\documentclass[11pt,a4paper]{article}
% preamble.tex — estilo común de todos los apuntes Froth.
% Cada apunte hace % preamble.tex — estilo común de todos los apuntes Froth.
% Cada apunte hace % preamble.tex — estilo común de todos los apuntes Froth.
% Cada apunte hace \input{preamble} justo después de \documentclass.
\usepackage[margin=2.4cm]{geometry}
\usepackage{xcolor}
\definecolor{litblue}{RGB}{31,79,121}
\definecolor{litgreen}{RGB}{27,94,32}
\definecolor{litorange}{RGB}{230,81,0}
\usepackage{parskip}
\usepackage{titlesec}
\titleformat{\section}{\Large\bfseries\color{litblue}}{\thesection}{0.6em}{}
\titleformat{\subsection}{\large\bfseries\color{litblue}}{\thesubsection}{0.6em}{}
\usepackage{tcolorbox}
\usepackage{fancyhdr}
\pagestyle{fancy}
\fancyhf{}
\fancyhead[L]{\small\color{litblue}Froth — Apuntes de ML}
\fancyhead[R]{\small\thepage}
\renewcommand{\headrulewidth}{0.4pt}
\usepackage[colorlinks=true,linkcolor=litblue,urlcolor=litblue]{hyperref}

% Cabecera de cada apunte: \cabecera{numero}{titulo}
\newcommand{\cabecera}[2]{%
  \begin{tcolorbox}[colback=litblue,coltext=white,colframe=litblue,arc=2mm]
    {\Large\bfseries Apunte #1 — #2}\\[3pt]
    {\small Proyecto Froth · aprender Machine Learning construyendo}
  \end{tcolorbox}\vspace{0.8em}}

% Cajas temáticas
\newtcolorbox{concepto}[1]{colback=blue!4,colframe=litblue,title={Concepto: #1},arc=1.5mm}
\newtcolorbox{practica}{colback=green!5,colframe=litgreen,title={Pruébalo tú (teclado en mano)},arc=1.5mm}
\newtcolorbox{ojo}{colback=orange!8,colframe=litorange,title={Ojo / error típico},arc=1.5mm}
\newtcolorbox{resumen}{colback=gray!8,colframe=gray!60,title={En una frase},arc=1.5mm}

\newcommand{\codigo}[1]{\texttt{#1}}
 justo después de \documentclass.
\usepackage[margin=2.4cm]{geometry}
\usepackage{xcolor}
\definecolor{litblue}{RGB}{31,79,121}
\definecolor{litgreen}{RGB}{27,94,32}
\definecolor{litorange}{RGB}{230,81,0}
\usepackage{parskip}
\usepackage{titlesec}
\titleformat{\section}{\Large\bfseries\color{litblue}}{\thesection}{0.6em}{}
\titleformat{\subsection}{\large\bfseries\color{litblue}}{\thesubsection}{0.6em}{}
\usepackage{tcolorbox}
\usepackage{fancyhdr}
\pagestyle{fancy}
\fancyhf{}
\fancyhead[L]{\small\color{litblue}Froth — Apuntes de ML}
\fancyhead[R]{\small\thepage}
\renewcommand{\headrulewidth}{0.4pt}
\usepackage[colorlinks=true,linkcolor=litblue,urlcolor=litblue]{hyperref}

% Cabecera de cada apunte: \cabecera{numero}{titulo}
\newcommand{\cabecera}[2]{%
  \begin{tcolorbox}[colback=litblue,coltext=white,colframe=litblue,arc=2mm]
    {\Large\bfseries Apunte #1 — #2}\\[3pt]
    {\small Proyecto Froth · aprender Machine Learning construyendo}
  \end{tcolorbox}\vspace{0.8em}}

% Cajas temáticas
\newtcolorbox{concepto}[1]{colback=blue!4,colframe=litblue,title={Concepto: #1},arc=1.5mm}
\newtcolorbox{practica}{colback=green!5,colframe=litgreen,title={Pruébalo tú (teclado en mano)},arc=1.5mm}
\newtcolorbox{ojo}{colback=orange!8,colframe=litorange,title={Ojo / error típico},arc=1.5mm}
\newtcolorbox{resumen}{colback=gray!8,colframe=gray!60,title={En una frase},arc=1.5mm}

\newcommand{\codigo}[1]{\texttt{#1}}
 justo después de \documentclass.
\usepackage[margin=2.4cm]{geometry}
\usepackage{xcolor}
\definecolor{litblue}{RGB}{31,79,121}
\definecolor{litgreen}{RGB}{27,94,32}
\definecolor{litorange}{RGB}{230,81,0}
\usepackage{parskip}
\usepackage{titlesec}
\titleformat{\section}{\Large\bfseries\color{litblue}}{\thesection}{0.6em}{}
\titleformat{\subsection}{\large\bfseries\color{litblue}}{\thesubsection}{0.6em}{}
\usepackage{tcolorbox}
\usepackage{fancyhdr}
\pagestyle{fancy}
\fancyhf{}
\fancyhead[L]{\small\color{litblue}Froth — Apuntes de ML}
\fancyhead[R]{\small\thepage}
\renewcommand{\headrulewidth}{0.4pt}
\usepackage[colorlinks=true,linkcolor=litblue,urlcolor=litblue]{hyperref}

% Cabecera de cada apunte: \cabecera{numero}{titulo}
\newcommand{\cabecera}[2]{%
  \begin{tcolorbox}[colback=litblue,coltext=white,colframe=litblue,arc=2mm]
    {\Large\bfseries Apunte #1 — #2}\\[3pt]
    {\small Proyecto Froth · aprender Machine Learning construyendo}
  \end{tcolorbox}\vspace{0.8em}}

% Cajas temáticas
\newtcolorbox{concepto}[1]{colback=blue!4,colframe=litblue,title={Concepto: #1},arc=1.5mm}
\newtcolorbox{practica}{colback=green!5,colframe=litgreen,title={Pruébalo tú (teclado en mano)},arc=1.5mm}
\newtcolorbox{ojo}{colback=orange!8,colframe=litorange,title={Ojo / error típico},arc=1.5mm}
\newtcolorbox{resumen}{colback=gray!8,colframe=gray!60,title={En una frase},arc=1.5mm}

\newcommand{\codigo}[1]{\texttt{#1}}

\begin{document}
\cabecera{45}{Zoom semántico: por qué 10.921 papers caben en pantalla y 1.772 no}

\section{Glosario en palabras simples}

\begin{itemize}
  \item \textbf{Renderizar}: dibujar en la pantalla. Cada punto, cada línea, cada letra
        que ves alguien tuvo que pintarla.
  \item \textbf{DOM}: la lista de objetos que el navegador guarda para cada cosa
        dibujada en una página. Si dibujas 10.000 círculos como objetos, el navegador
        guarda 10.000 fichas y las revisa continuamente.
  \item \textbf{SVG}: forma de dibujar donde CADA figura es un objeto del DOM. Cómoda
        (puedes darle clic a cada una) pero cara: miles de objetos ahogan al navegador.
  \item \textbf{WebGL}: forma de dibujar que le pasa los puntos a la \textbf{tarjeta
        gráfica} (la GPU). La GPU pinta millones de puntos a la vez porque para eso está
        hecha. No guarda una ficha por punto.
  \item \textbf{GPU}: el chip que dibuja. Tu portátil tiene una RTX 5060. Hace una cosa
        muy bien: la misma operación sobre miles de datos \emph{en paralelo}.
  \item \textbf{Hilo principal} (\emph{main thread}): el único carril por donde el
        navegador atiende TODO: tus clics, la rueda del ratón, los cálculos, el dibujo.
        Si algo lo ocupa, la página deja de responder aunque no esté ``rota''.
  \item \textbf{Fotograma} (\emph{frame}): una imagen de la pantalla. Para que algo se
        sienta fluido hacen falta unos 60 por segundo, o sea $\approx$16 ms para cada uno.
  \item \textbf{Física de grafo}: simulación donde los nodos se repelen y las aristas
        tiran como muelles, hasta que el dibujo se ``acomoda'' solo.
  \item \textbf{Percentil}: si ordenas 201 números de menor a mayor, el percentil 90 es
        el valor que deja por debajo al 90\% de ellos. Sirve para decir ``los grandes''
        sin inventarse un número.
  \item \textbf{Centroide}: el punto medio de un grupo. Donde colocamos su nombre.
\end{itemize}

\section{El problema, tal como apareció}

Subimos el tope de papers de la vista \emph{Network} de 400 a 2.500 y quedaron 1.772
nodos en pantalla. Yo escribí ``no se rompe nada''. Tu pantallazo dijo otra cosa: una
bola de pelo congelada, sin zoom ni desplazamiento.

Y sin embargo, en la misma sesión, la vista \emph{Atlas} dibujó \textbf{10.921 papers} y
respondía. Casi seis veces más puntos, y fluida. Ese contraste es el apunte entero.

\begin{concepto}{Dibujar no es simular}
La Network usa \textbf{física}: en cada fotograma recalcula las fuerzas entre nodos y,
cuando mueves el ratón, comprueba nodo por nodo cuál está debajo del cursor. Con 1.772
nodos eso llena los 16 ms del fotograma y no queda tiempo para atender la rueda. El zoom
no estaba roto: \textbf{el carril estaba ocupado}.

El Atlas no simula nada. Las posiciones ya están calculadas de antes (UMAP, apunte 08) y
guardadas. Solo las manda a la GPU y esta las pinta. Por eso 10.921 le cuestan menos que
1.772 con física.
\end{concepto}

\begin{ojo}
\textbf{Mi error, con todas las letras.} Cuando dije ``no se rompe nada'' yo había medido dos cosas: cuánto tarda Python en
\emph{generar} el HTML, y cuánto tarda vis.js en \emph{asentar} el dibujo. Las dos
baratas. \textbf{No medí la interacción}, que es justo lo que se rompió.

La lección, que vale para todo el proyecto: \textbf{que algo sea rápido de PRODUCIR no
dice nada sobre si es fluido de USAR}. Son dos medidas distintas y hay que tomar las dos.
\end{ojo}

\section{La otra mitad: 201 nombres no caben}

Arreglado el motor de dibujo, aparece el problema humano. El Atlas pintaba los puntos con
un color por subtema y nada más: bonito, y mudo. Para saber qué era cada mancha había que
irse a la leyenda lateral y cruzar colores a ojo.

La solución obvia (poner el nombre de cada subtema sobre su centroide) no funciona:
en el corpus de LIBS/FTIR hay \textbf{201 subtemas}. Doscientos un rótulos a la vez son
una pared de texto, no un mapa.

\subsection{Mini-ejemplo de juguete (antes de la teoría)}

Imagina un mapa de solo 5 territorios, con este tamaño en papers:

\begin{center}
\begin{tabular}{lccccc}
\hline
Territorio & A & B & C & D & E \\
Papers     & 100 & 60 & 30 & 12 & 8 \\
\hline
\end{tabular}
\end{center}

Ahora piensa en Google Maps. Muy lejos ves ``España''. Te acercas y aparece ``Madrid''.
Más cerca, ``Chamberí''. Más cerca todavía, el nombre de la calle. \textbf{Lo que se
muestra cambia con la distancia}, no solo el tamaño de las letras. Eso es zoom semántico.

Aplicado a nuestros 5 territorios: desde lejos solo A; te acercas y aparecen A y B;
más cerca A, B, C; encima del todo, los cinco. La pregunta interesante es
\textbf{dónde pones los cortes}, y ahí es donde no vale inventarse un número.

\begin{concepto}{Zoom semántico (Shneiderman, 1996)}
La regla clásica de visualización dice: \emph{``vista general primero, luego zoom y
filtro, y el detalle solo cuando se pide''}. Traducido a Froth: primero unos pocos
territorios grandes con nombre, y los subtemas pequeños solo cuando te acercas a su zona.
Nunca los 201 de golpe.
\end{concepto}

\section{Los cortes salen de TUS datos (la regla de siempre)}

Un umbral inventado (``muestra los que tengan más de 50 papers'') funcionaría en un
corpus y sería absurdo en otro. Así que los cortes se calculan con los
\textbf{percentiles del tamaño de los subtemas de ese corpus concreto}: percentil 90
(los grandes de verdad), 75, 50, y 0 (todos).

Dos corpus reales tuyos, con los números que salieron:

\begin{center}
\begin{tabular}{lrrl}
\hline
Corpus & Papers & Subtemas & Cortes calculados (p90/p75/p50/p0) \\
\hline
Beauvoir (tu tesis)  & 3.867  & 96  & 54 / 35 / 19 / 7 \\
LIBS + FTIR          & 10.921 & 201 & 76 / 37 / 20 / 7 \\
Economía circular    & 4.946  & 69  & 142 / 69 / 33 / 7 \\
\hline
\end{tabular}
\end{center}

Fíjate en la tercera fila: economía circular tiene MENOS papers que LIBS pero su corte de
``satélite'' es casi el doble (142 frente a 76), porque sus subtemas son más gordos y
menos numerosos. Un número fijo habría dejado ese mapa vacío o saturado. El corpus decide.

Lo que ves al acercarte, medido en el navegador sobre los 10.921 papers:

\begin{center}
\begin{tabular}{llrr}
\hline
Nivel & Metáfora & Corte & Nombres en pantalla \\
\hline
1 & Satélite & $\geq$ 76 & 21 \\
2 & Regional & $\geq$ 37 & 52 \\
3 & Local    & $\geq$ 20 & 106 \\
4 & Calle    & $\geq$ 7  & 201 \\
\hline
\end{tabular}
\end{center}

Y todavía hay una segunda red de seguridad: aunque en el nivel 4 se manden los 201
rótulos, deck.gl aplica un \textbf{filtro de colisión} que descarta los que se pisarían,
dejando ganar al territorio más grande. Los 201 solo pueden aparecer juntos si de verdad
caben.

\begin{practica}
Abre el Atlas de un corpus grande y haz esto en orden:
\begin{enumerate}
  \item Sin tocar nada, cuenta los nombres que ves. Deberían ser pocos y todos de
        territorios gordos.
  \item Acerca la rueda dos o tres pasos sobre una zona. Verás aparecer nombres nuevos,
        y solo en la zona donde estás mirando.
  \item Pulsa un subtema en la leyenda lateral: la cámara vuela hasta él.
  \item Pulsa \textbf{``Whole map''} para volver a la vista general.
\end{enumerate}
Si en el paso 1 ves una pared de texto, el cálculo de percentiles no se está aplicando:
avísame y lo miramos.
\end{practica}

\section{El tropiezo que enseña más que el éxito}

El plan original para ``clic en un subtema'' era distinto y más ambicioso: recalcular un
UMAP \textbf{solo con los papers de ese subtema}, para revelar su estructura interna, que
el UMAP global aplana. El plan decía ``unos 2 segundos, y se puede cachear''.

Antes de construirlo, lo medí: \textbf{29,2 segundos} para 500 papers, con la función
real (\texttt{cluster.reduce\_to\_2d}). No 2 segundos. Casi treinta.

Con 96 subtemas en tu corpus de tesis, explorar el mapa habría sido esperar medio minuto
por cada clic. Inaceptable. Así que el clic hace algo mucho más humilde: \textbf{vuela la
cámara a las coordenadas que ya existen}. Cero cálculo, cero espera.

\begin{ojo}
\textbf{La cifra del plan era mía, y era falsa.} Ese ``$\sim$2s'' no venía de ninguna medición: lo escribí yo por intuición y quedó en el
plan como si fuera un dato. Si no lo llego a medir antes de programar, habríamos
construido una función inservible y lo habrías descubierto tú, usándola.

Regla: \textbf{una cifra en un plan no es un dato hasta que alguien la mide}. Y las
mediciones se hacen ANTES de construir, no después.
\end{ojo}

Se pierde algo real, y conviene decirlo: la sub-estructura interna de cada subtema sigue
sin verse. Queda anotado como posible mejora futura (calcularla una sola vez, en el
momento de procesar el corpus, no al hacer clic).

\section{La división del trabajo que salió de aquí}

En vez de intentar que una vista lo hiciera todo, cada una hace lo que sabe hacer:

\begin{center}
\begin{tabular}{lll}
\hline
 & \textbf{Network} & \textbf{Atlas} \\
\hline
Motor          & vis.js con física   & deck.gl / WebGL (GPU) \\
Tope de papers & 400                 & el corpus entero \\
Qué muestra    & quién se parece a quién & el territorio completo \\
Para qué sirve & leer el núcleo más citado & orientarte y elegir dónde mirar \\
\hline
\end{tabular}
\end{center}

Bajar la Network de 2.500 a 400 \textbf{no es una regresión}: es reconocer que la física
sirve para ver relaciones entre pocos nodos, y estorba para ver un continente. Si 400 te
sabe a poco, el número está en la configuración y se sube sin tocar código; pero la
fluidez solo la puede juzgar tu máquina, no la mía.

\begin{resumen}
Dibujar en la GPU (WebGL) escala; simular física por fotograma, no. Y cuando hay más
nombres que sitio, se muestran por niveles de zoom, con los cortes calculados a partir de
los percentiles de TU corpus, no de un número inventado.
\end{resumen}

\section{Para saber más}

\begin{itemize}
  \item Shneiderman, B. (1996), \emph{The Eyes Have It}: el origen del ``overview first,
        zoom and filter, then details-on-demand''.
  \item \texttt{nightingaledvs.com/how-to-visualize-a-graph-with-a-million-nodes}
  \item \texttt{cambridge-intelligence.com/visualizing-graphs-webgl}
  \item El código: \texttt{froth/atlas.py} (funciones \texttt{minCountForZoom},
        \texttt{makeLabels} y \texttt{flyTo}).
\end{itemize}

\end{document}
